% PAGES: 131-132

% Continuation of sec04_c1_2.tex (page 130 / folio 114): this text concludes
% the \item ("If lambda = 1, any corresponding eigenvector...") for matrix A,
% inside the \begin{itemize} and the manual "Examples:" block opened there.
% Do not reopen either.
Thus, any eigenvector corresponding to the eigenvalue $\lambda = 1$ is of the form
\[
v \;=\; \begin{pmatrix}v_1\\ -0.5v_1\\ 3v_1\end{pmatrix} \;=\;
0.5v_1\begin{pmatrix}2\\-1\\6\end{pmatrix}.
\]
We conclude that the eigenvalue $\lambda = 1$ has one eigenvector, e.g.,
$\begin{pmatrix}2\\-1\\6\end{pmatrix}$.
\end{itemize}

(ii) Let
\[
B \;=\; \begin{pmatrix} 2 & 2 & -1\\ -7 & -4 & 2.5\\ 2 & 4 & -1\end{pmatrix}.
\]
The characteristic polynomial of $B$ is
\[
P_B(t) \;=\; \det(tI-B) \;=\;
\begin{vmatrix} t-2 & -2 & 1\\ 7 & t+4 & -2.5\\ -2 & -4 & t+1\end{vmatrix}.
\]
Using formula (10.8), we find that
\[
P_B(t) \;=\; t^3 + 3t^2 - 4 \;=\; (t-1)(t+2)^2.
\]
The roots of the characteristic polynomial $P_B(t)$ are $\lambda_1 = \lambda_2 = -2$ and
$\lambda_3 = 1$.

\begin{itemize}
\item If $\lambda = -2$, any corresponding eigenvector $v \neq 0$ is a solution to
$Bv = -2v$, which can be written as
\[
\begin{cases}
2v_1+2v_2-v_3 = -2v_1\\
-7v_1-4v_2+2.5v_3 = -2v_2\\
2v_1+4v_2-v_3 = -2v_3
\end{cases}
\quad\Longleftrightarrow\quad
\begin{cases}
v_3 = 4v_1+2v_2\\
2.5v_3 = 7v_1+2v_2\\
v_3 = -2v_1-4v_2
\end{cases}
\]
\[
\Longleftrightarrow\quad
\begin{cases}
v_3 = 4v_1+2v_2\\
5v_3 = 14v_1+4v_2\\
0 = 6v_1+6v_2
\end{cases}
\quad\Longleftrightarrow\quad
\begin{cases}
v_3 = 2v_1\\
v_3 = 2v_1\\
v_2 = -v_1
\end{cases}
\]
Thus, any eigenvector corresponding to the eigenvalue $\lambda = -2$ is of the form
\[
v \;=\; \begin{pmatrix}v_1\\ -v_1\\ 2v_1\end{pmatrix} \;=\;
v_1\begin{pmatrix}1\\-1\\2\end{pmatrix}
\]
and we conclude that the eigenvalue $\lambda = -2$ has one eigenvector, e.g.,
$\begin{pmatrix}1\\-1\\2\end{pmatrix}$.
\item If $\lambda = 1$, any corresponding eigenvector $v \neq 0$ is a solution to
$Bv = v$, which can be written as
\[
\begin{cases}
2v_1+2v_2-v_3 = v_1\\
-7v_1-4v_2+2.5v_3 = v_2\\
2v_1+4v_2-v_3 = v_3
\end{cases}
\quad\Longleftrightarrow\quad
\begin{cases}
v_3 = v_1+2v_2\\
5v_3 = 14v_1+10v_2\\
v_3 = v_1+2v_2
\end{cases}
\]
\[
\Longleftrightarrow\quad
\begin{cases}
v_3 = v_1+2v_2\\
0 = 9v_1\\
v_3 = v_1+2v_2
\end{cases}
\quad\Longleftrightarrow\quad
\begin{cases}
v_3 = 2v_2\\
v_1 = 0\\
v_3 = 2v_2
\end{cases}
\]
Thus, any eigenvector corresponding to the eigenvalue $\lambda = 1$ is of the form
\[
v \;=\; \begin{pmatrix}0\\ v_2\\ 2v_2\end{pmatrix} \;=\; v_2\begin{pmatrix}0\\1\\2\end{pmatrix}.
\]
We conclude that the eigenvalue $\lambda = 1$ has one eigenvector, e.g.,
$\begin{pmatrix}0\\1\\2\end{pmatrix}$.
\end{itemize}
\hfill$\square$\par\medskip

\begin{examplex}
The characteristic polynomial of the $2 \times 2$ matrix $A = \begin{pmatrix} a & b\\ c &
d\end{pmatrix}$ is\footnote{Note that (4.7) corresponds to (4.19), since the trace and
determinant of the $2 \times 2$ matrix $A$ are $\tr(A) = a+d$ and $\det(A) = ad-bc$,
respectively.}
\begin{equation}
P_A(t) \;=\; t^2 - (a+d)t + ad - bc,
\end{equation}
since, from definition (4.5) and the formula (10.6) for the determinant of a
$2 \times 2$ matrix, it follows that
\begin{align*}
P_A(t) \;=\; \det(tI-A) \;=\; \det\begin{pmatrix} t-a & -b\\ -c & t-d\end{pmatrix}
&\;=\; (t-a)(t-d) - (-b)(-c)\\
&\;=\; t^2 - (a+d)t + ad - bc.
\end{align*}
\end{examplex}

% The "Examples:" (plural) discussion below -- diagonal, lower triangular, and
% upper triangular matrices -- is written manually for the same reason as the
% previous "Examples:" block (the source label is plural; \examplex hard-codes
% the singular label). It is NOT closed in this file: part (iii), the upper
% triangular matrix $U$, and the closing "\hfill$\square$\par\medskip" are in
% sec04_c1_4.tex (page 133 / folio 117). sec04_c1_4.tex must not print a
% second "Examples:" label for this block.
\par\medskip\noindent\textit{Examples:}\ (i) Let $D = \diag(d_i)_{i=1:n}$ be a diagonal
matrix of size $n$. The characteristic polynomial of $D$ is
\[
P_D(t) \;=\; \det(tI-D) \;=\; \prod_{i=1}^n (t-d_i);
\]
cf. (10.1). The roots of the polynomial $P_D(t)$, i.e., the solutions to $P_D(t) = 0$,
are $d_i$, $i = 1:n$. We conclude that the eigenvalues of $D$ are its diagonal entries.
Note that the corresponding eigenvector for the eigenvalue $d_i$ is $e_i$, the $i$-th
column of the identity matrix, i.e.,
\[
De_i \;=\; d_ie_i, \quad \forall\, i = 1:n.
\]
Thus, a diagonal matrix of size $n$ has $n$ eigenvectors.

(ii) Let $L$ be a lower triangular matrix of size $n$. The characteristic polynomial of
$L$ is
\[
P_L(t) \;=\; \det(tI-L) \;=\; \prod_{i=1}^n (t-L(i,i));
\]
cf. (10.3). The roots of $P_L(t)$, i.e., the solutions to $P_L(t) = 0$, are $L(i,i)$,
$i = 1:n$. We conclude that the eigenvalues of $L$ are its diagonal entries $L(i,i)$,
$i = 1:n$. Note that the eigenvector corresponding to the eigenvalue $L(n,n)$ is $e_n$,
the $n$-th column of the identity matrix. However, the eigenvector corresponding to the
eigenvalue $L(i,i)$, with $1 \le i \le (n-1)$, is not necessarily $e_i$, the $i$-th
column of the identity matrix.
% CONTINUES in sec04_c1_4.tex (page 133 / folio 117): part (iii), the upper
% triangular matrix U, picks up here. The manual "Examples:" block above is
% left OPEN and is closed there with "\hfill$\square$\par\medskip".
